\documentclass[12pt, a4paper]{article}

% --- Packages ---
\usepackage[margin=1in]{geometry}
\usepackage{amsmath, amssymb, amsthm}
\usepackage{setspace}
\usepackage{hyperref}
\usepackage{natbib}
\usepackage{booktabs}
\usepackage{enumitem}
\usepackage{titlesec}

% --- Theorem environments ---
\newtheorem{proposition}{Proposition}
\newtheorem{lemma}{Lemma}
\newtheorem{corollary}{Corollary}
\newtheorem{remark}{Remark}
\newtheorem{definition}{Definition}

% --- Formatting ---
\onehalfspacing
\titleformat{\section}{\large\bfseries}{\thesection.}{0.5em}{}
\titleformat{\subsection}{\normalsize\bfseries}{\thesubsection.}{0.5em}{}

% --- Title ---
\title{\textbf{Referee Report: \\ ``Consumer Information and the Limits to Competition''} \\[6pt]
\large Mark Armstrong \& Jidong Zhou \\
\textit{American Economic Review}, 2022, 112(2), 534--577}
\author{Jordi Torres Vallverd\'{u} \\
\small Toulouse School of Economics --- MRes IO Theory, 2025--2026}
\date{\today}

\begin{document}

\maketitle
\thispagestyle{empty}

% =============================================
% PART 1: SUMMARY (2-4 pages)
% =============================================
\section{Summary}

\subsection{Motivation}

In many markets, firms supply differentiated products and compete in prices, while consumers possess incomplete information about how well each product matches their preferences. A central question for platforms (e.g., Amazon, Expedia, insurance comparison websites) and regulators is how to design the information environment that consumers face.

The key tension is immediate: more transparency improves product-consumer matching, but it also softens price competition by granting firms greater market power. Less transparency forces fierce price competition --- products are perceived as near-perfect substitutes --- but consumers may end up purchasing a poorly matched product.

Armstrong and Zhou formalize this trade-off. They ask: within the class of all feasible signal structures, which one maximizes industry profit, and which one maximizes consumer surplus? Their answer reveals that the interests of firms and consumers are sharply opposed in how information should be structured, yet both sides agree that information should be \emph{symmetric} across firms.

\subsection{Model Essentials}

Two symmetric firms ($i = 1,2$) supply differentiated products at zero cost. A risk-neutral consumer values product $i$ at $v_i \geq 0$; only her relative preference $x \equiv v_1 - v_2$ matters. Let $F$ denote the prior CDF of $x$, symmetric on $[-1,1]$, with density $f$. Two key parameters summarize the prior:
\begin{itemize}[nosep]
    \item $\mu \equiv E[v_i]$: expected value of a random product.
    \item $\delta \equiv E_F[\max\{x,0\}] = \int_{-1}^{0} F(x)\,dx$: expected match-quality gain from buying the preferred product over a random one.
\end{itemize}

Before buying, the consumer observes a private signal $s$ from a signal structure $\{\sigma(s|x), S\}$ chosen by a designer. She updates to a posterior; since she is risk-neutral, only the posterior mean $E[x \mid s]$ matters. Let $G$ denote the CDF of this posterior mean --- the \textbf{posterior distribution}. Firms observe $G$ but not $s$, and set prices simultaneously (Bertrand--Nash). The consumer buys from firm~1 if $E[x \mid s] > p_1 - p_2$.

The only constraint on $G$ is Bayesian consistency: $G$ must be a mean-preserving contraction (MPC) of $F$, i.e.,
$$\int_{-1}^{x} G(\tilde{x})\,d\tilde{x} \;\leq\; \int_{-1}^{x} F(\tilde{x})\,d\tilde{x} \qquad \forall\, x \in [-1,1].$$
Any such $G$ can be generated by some signal structure. So instead of optimizing over signals, the authors optimize directly over posteriors --- a major simplification.

A posterior $G$ can sustain symmetric equilibrium price $p$ if and only if $G$ lies between two bounds derived from each firm's no-deviation condition:
$$L_p(x) \;\leq\; G(x) \;\leq\; U_p(x), \qquad U_p(x) = \frac{p}{2(p-x)}, \quad L_p(x) = 1 - U_p(-x).$$
Higher $p$ flattens both bounds, requiring fewer near-indifferent consumers. Consumer surplus under a symmetric posterior with price $p$ is $CS = \mu + \int_{-1}^{0} G(x)\,dx - p$.

\subsection{The Core Trade-Off}

Two polar cases make the tension concrete. Under \textbf{full disclosure} ($G = F$), the consumer always buys her preferred product (max welfare), but firms enjoy market power and charge high prices. Under \textbf{no disclosure} ($G$ degenerate at zero), products look identical, Bertrand competition drives prices to zero, but the consumer buys randomly and may get the wrong product. Surprisingly, the consumer can be \emph{better off} with no information: the price drop more than compensates the mismatch.

The paper's contribution is to characterize the \emph{entire} spectrum between these extremes: which posteriors $G$ maximize firm profit, which maximize consumer surplus, and what combinations of profit and surplus are feasible.

\subsection{Key Insight: Symmetric Signals Dominate}

Before solving the optimal design, the paper establishes a powerful restriction: for \emph{any} asymmetric signal (one that favors a firm), there exists a symmetric signal that yields higher profit for \emph{each} firm and higher consumer surplus. Biased recommendations trigger aggressive undercutting by the disfavored firm, hurting everyone. This means we can restrict attention to symmetric posteriors and symmetric prices throughout.

\subsection{Main Results}

\textbf{Firm-optimal signal.} The firm-optimal posterior sits on the \emph{lower} bound of the feasibility corridor, yielding a U-shaped density: consumers are pushed toward the extremes, few are near-indifferent, and neither firm can profitably undercut. The resulting price is
$$p^* = \frac{2\delta}{1 - \log 2} \approx 6.5\,\delta,$$
where $\delta \equiv E_F[\max\{x,0\}]$ measures underlying product differentiation. There is \textbf{no mismatch}: the consumer buys her preferred product, and total welfare is maximized.

\textbf{Consumer-optimal signal.} The consumer-optimal posterior rides the \emph{upper} bound: mass is concentrated near $x = 0$, creating many swing consumers who make firms compete fiercely. The price drops to $p^{**} \approx 0.05\, p_F$. The cost is \textbf{mismatch} --- some consumers buy the wrong product --- but the price effect dominates.

\textbf{The limits to competition.} Any profit $p \in [0, p^*]$ is feasible. Competition makes it \emph{impossible} to have both low profit and low consumer surplus simultaneously --- unlike the monopoly setting of Roesler and Szentes (2017). The feasible set has a concave frontier, and full disclosure is optimal only for unweighted total welfare.


% =============================================
% PART 2: CRITIQUE (1-2 pages)
% =============================================
\section{Critique}

While the paper makes an important and technically elegant contribution, several aspects of the modeling framework deserve scrutiny.

\subsection{Scalar Heterogeneity and the Duopoly Restriction}

The paper reduces the consumer's two-dimensional preference $(v_1, v_2)$ to a scalar relative valuation $x = v_1 - v_2$. This is natural in the duopoly case, but it means that the \emph{level} of consumer willingness-to-pay (captured by $V = v_1 + v_2$) is effectively assumed away. In practice, both the level and the relative preference matter: a consumer who values both products highly behaves differently from one who values both poorly, even if their relative preference $x$ is identical. With more than two firms, preferences are inherently multidimensional, and the paper's approach does not extend. The authors study three simple policies (``top product,'' ``top two,'' and full disclosure) for $n > 2$, but the full characterization of optimal signal structures remains open.

\subsection{Pure-Strategy Restriction}

The focus on signal structures that induce pure-strategy pricing equilibria is motivated by tractability. Proposition~4 shows that allowing mixed strategies improves consumer surplus by at most 1.6\%, which is reassuring. However, the firm-optimal signal under mixed strategies remains uncharacterized. More importantly, for some prior distributions and signal structures, pure-strategy equilibria may not exist, and the paper's framework provides no guidance for these cases. Understanding whether the qualitative insights --- symmetric dominance, the shape of the feasible frontier --- survive under mixed strategies would strengthen the paper's conclusions.

\subsection{Exogeneity of the Signal Structure}

The entire analysis assumes the signal structure is chosen exogenously by a designer and committed to before firms set prices. This ``centralized design'' assumption is strong. In many real markets, information is generated through a combination of firm advertising, consumer search, platform algorithms, and third-party reviews.
\begin{itemize}[nosep]
    \item If firms can influence the signal structure (e.g., through advertising), the optimal centralized policy may not be implementable.
    \item If consumers can choose how much information to acquire (endogenous search), the designer's ability to control the posterior is limited.
    \item The timing assumption is crucial. If firms could respond to the signal structure by adjusting product characteristics or quality, the welfare conclusions could reverse.
\end{itemize}

\subsection{Symmetry of the Prior}

The paper assumes firms are ex ante symmetric: the prior over $(v_1, v_2)$ is symmetric between the two products. Many real-world markets feature vertical differentiation (one product is known to be better on average) or brand effects. The authors suggest (Section~IV) that extending to asymmetric priors is interesting, but they do not explore whether the optimality of symmetric signals (Lemma~2) survives. With an asymmetric prior, biased signals that further favor the better firm might no longer be dominated, and the structure of the feasible set could change qualitatively.

\subsection{Welfare Metric and Dynamic Considerations}

The consumer-optimal signal structure deliberately introduces mismatch: some consumers buy their less-preferred product. The paper treats this as a pure static efficiency loss, measured by $\delta - \int_{-1}^{0} G(x)\,dx$. However, in practice, product mismatch may generate dynamic costs --- consumer regret, returns, negative reviews, reduced repeat purchases --- that are not captured in the one-shot model. If these costs are substantial, the consumer-optimal policy may be less attractive than the static analysis suggests, and a regulator might optimally choose a signal structure that lies between the firm-optimal and consumer-optimal extremes.


% =============================================
% PART 3: EXTENSION (2-4 pages)
% =============================================
\section{Extension: Information Designer as a Strategic Agent}

\subsection{Motivation}

Armstrong and Zhou explicitly call for embedding the information designer as a strategic agent (Section~IV, p.~562). In their framework, signal structures are chosen to maximize firm profit or consumer surplus, but no agent actually performs this optimization. A natural candidate is a \textbf{platform} --- think of Amazon or Expedia --- that controls how much product information consumers see and monetizes through fees to firms. The question is: \emph{what signal structure would such a platform choose, and how does its answer depend on how it charges?}

I sketch below how the paper's existing tools --- the bounds $L_p, U_p$, the MPC constraint, and the feasible set (Proposition~3) --- can be used to reason about what the platform would find optimal under two canonical fee structures. The arguments are informal and intended as a direction for further investigation, not as formal results.

\subsection{Setup}

We add a platform to the Armstrong--Zhou duopoly. The platform commits to a signal structure (posterior $G$) and a fee \emph{before} firms set prices. Everything else is unchanged: two symmetric firms, zero cost, prior $F$ on $[-1,1]$, consumer always buys one product. We consider two fee instruments:

\begin{itemize}[nosep]
    \item \textbf{Per-transaction fee $\tau$}: the platform receives $\tau$ on every sale.
    \item \textbf{Lump-sum fee $\phi$}: each firm pays $\phi$ upfront to be listed.
\end{itemize}

\subsection{A Useful Observation: Per-Unit Fees and Markups}

With per-unit fee $\tau$, each firm's effective cost becomes $\tau$. Define the markup $\tilde{p}_i \equiv p_i - \tau$. The consumer's choice depends only on the price \emph{difference} $p_1 - p_2 = \tilde{p}_1 - \tilde{p}_2$, so the demand system in markups is identical to the original game. Each firm maximizes $\tilde{p}_i \cdot q_i(\tilde{p}_1, \tilde{p}_2; G)$ --- the same problem as with zero cost. This suggests that the equilibrium markup $\tilde{p} = p(G)$ should be unchanged: the fee is fully passed through, the consumer pays $p(G) + \tau$, and firm profits $\pi = p(G)/2$ are \textbf{invariant} to $\tau$.

(This pass-through argument relies on the structure of Bertrand competition with inelastic total demand. A full verification would check that the no-deviation bounds $L_p, U_p$ are unchanged in markups, which follows from the price difference being the only relevant variable for demand.)

\subsection{Sketch: What Would the Platform Choose?}

\subsubsection{Case 1: Per-unit fees}

The platform earns $\tau$ per sale (total revenue $= \tau$, since one unit is always sold). It wants $\tau$ as large as possible. The binding constraint is consumer participation: the consumer should want to buy. Her surplus under posterior $G$ and fee $\tau$ is:
$$CS = \underbrace{\mu + \int_{-1}^{0} G(x)\,dx}_{\text{gross utility}} - \underbrace{\big(p(G) + \tau\big)}_{\text{total price}}.$$
Setting $CS = 0$ gives $\tau = \mu + \int_{-1}^{0} G(x)\,dx - p(G)$. The platform's revenue equals whatever consumer surplus the signal $G$ would generate in the original game --- and then extracts all of it via $\tau$.

So the platform's problem should reduce to: \emph{choose $G$ to maximize consumer surplus in the Armstrong--Zhou game}. If this reasoning is correct, the platform would want the consumer-optimal signal from Proposition~2: $G = \min\{F, U_{p^{**}}\}$ with $p^{**} \approx 0.05\, p_F$. It would then set $\tau$ to extract the surplus entirely.

\textbf{Conjectured outcome:} the consumer-optimal signal is implemented, but consumers get \emph{zero} surplus ($\tau$ takes it all). Firms earn small but positive markups. Mismatch occurs, total welfare is below $\mu + \delta$.

\subsubsection{Case 2: Lump-sum fees}

The platform charges $\phi$ per firm upfront. Since lump-sum fees do not enter marginal cost, they should not affect pricing or consumer surplus. Each firm participates iff $\pi(G) = p(G)/2 \geq \phi$. The platform sets $\phi = p(G)/2$ and earns $2\phi = p(G)$.

So the platform's problem should reduce to: \emph{choose $G$ to maximize industry profit}. If so, the platform would want the firm-optimal signal from Proposition~1: $G = L_{p^*}$ with $p^* = 2\delta/(1 - \log 2) \approx 6.5\,\delta$, and extract all firm profits via $\phi$.

\textbf{Conjectured outcome:} the firm-optimal signal is implemented. No mismatch, total welfare maximized at $\mu + \delta$. Consumers retain surplus $\mu + \delta - p^* > 0$ (since the platform cannot extract from them with firm-side fees).

\subsection{Comparison}

\begin{center}
\renewcommand{\arraystretch}{1.4}
\begin{tabular}{l c c}
\toprule
 & \textbf{Per-unit fee $\tau$} & \textbf{Lump-sum fee $\phi$} \\
\midrule
Platform maximizes & Consumer surplus & Industry profit \\
Signal chosen & Consumer-optimal ($G^{**}$) & Firm-optimal ($L_{p^*}$) \\
Mismatch? & Yes & No \\
Total welfare & $< \mu + \delta$ & $= \mu + \delta$ (max) \\
Consumer surplus & $0$ (extracted via $\tau$) & $\mu + \delta - p^* > 0$ \\
Firm profit (net of fee) & $p^{**}/2 > 0$ & $0$ (extracted via $\phi$) \\
\bottomrule
\end{tabular}
\end{center}

\subsection{Discussion and Tentative Takeaways}

If the reasoning above is correct, several interesting points emerge.

\textbf{The fee structure shapes information design.} The platform designs information to maximize the surplus it can capture. Per-unit fees let it capture consumer surplus $\Rightarrow$ consumer-optimal signal. Lump-sum fees let it capture firm profits $\Rightarrow$ firm-optimal signal. The business model would not be a distributional detail --- it would determine the entire information environment.

\textbf{A possible paradox: lump-sum fees may be better for consumers.} Under lump-sum fees, the platform extracts from firms, not consumers. So it designs information to maximize firm profits (which it takes), and consumers keep the residual $\mu + \delta - p^*$. Under per-unit fees, the platform targets consumers directly and takes everything. Consumers end up with nothing, and mismatch reduces total welfare. The more ``consumer-friendly'' fee structure (commissions) could actually lead to worse consumer outcomes. This would be worth verifying formally.

\textbf{Regulatory angle.} If this logic holds, a regulator who cannot directly control information design might instead regulate the platform's fee structure. Mandating lump-sum fees would indirectly implement the welfare-maximizing signal.

\textbf{What I have not checked.} Several things need verification: (i) whether the pass-through argument holds exactly or only approximately under the bounds structure; (ii) whether the consumer participation constraint is the right way to close the per-unit fee case, or whether other constraints bind first; (iii) what happens with competing platforms, two-sided fees, or if the platform cannot commit to $G$ before firms price. These are natural next steps.

% =============================================
% REFERENCES
% =============================================
\section*{References}

\begin{itemize}[nosep, leftmargin=*]
    \item Armstrong, M. and J. Zhou (2022). ``Consumer Information and the Limits to Competition.'' \textit{American Economic Review}, 112(2), 534--577.
    \item Bergemann, D., B. Brooks, and S. Morris (2015). ``The Limits of Price Discrimination.'' \textit{American Economic Review}, 105(3), 921--957.
    \item Ivanov, M. (2013). ``Information Revelation in Competitive Markets.'' \textit{Economic Theory}, 52(1), 337--365.
    \item Johnson, J. and D. Myatt (2006). ``On the Simple Economics of Advertising, Marketing, and Product Design.'' \textit{American Economic Review}, 96(3), 756--784.
    \item Kamenica, E. and M. Gentzkow (2011). ``Bayesian Persuasion.'' \textit{American Economic Review}, 101(6), 2590--2615.
    \item Rochet, J.-C. and J. Tirole (2003). ``Platform Competition in Two-Sided Markets.'' \textit{Journal of the European Economic Association}, 1(4), 990--1029.
    \item Roesler, A.-K. and B. Szentes (2017). ``Buyer-Optimal Learning and Monopoly Pricing.'' \textit{American Economic Review}, 107(7), 2072--2080.
\end{itemize}

\end{document}
